% 柯伊伯带（综述）
% license CCBYSA3
% type Wiki

本文根据 CC-BY-SA 协议转载翻译自维基百科\href{https://en.wikipedia.org/wiki/Kuiper_belt}{相关文章}。

柯伊伯带（/ˈkaɪpər/ ⓘ）$^{[2]}$ 是位于太阳系外部的一种环星盘结构，其范围从海王星轨道处约 30 个天文单位（AU）一直延伸到距离太阳约 50 AU 的区域。$^{[3]}$ 它在结构上类似于小行星带，但规模要大得多——宽度约为其 20 倍，总质量约为其 20–200 倍。$^{[4][5]}$ 与小行星带相似，柯伊伯带主要由太阳系形成早期遗留下来的小天体或残骸组成。不同之处在于，许多小行星主要由岩石和金属构成，而大多数柯伊伯带天体则主要由冻结的挥发性物质（通常称为“冰”）构成，例如甲烷、氨和水冰。柯伊伯带是天文学界普遍认可的多颗矮行星的家园，包括 Orcus、冥王星（Pluto）$^{[6]}$、妊神星（Haumea）$^{[7]}$、夸奥瓦（Quaoar）以及鸟神星（Makemake）$^{[8]}$。太阳系中的一些卫星，例如海王星的海卫一（Triton）和土星的菲比（Phoebe），也可能起源于这一地区。$^{[9][10]}$

柯伊伯带以荷兰天文学家杰拉德·柯伊伯（Gerard Kuiper）的名字命名，他在 1951 年提出了这种带状结构存在的一种设想。$^{[11]}$ 在他之前和之后，也有研究者提出过类似的假说，例如 20 世纪 30 年代的肯尼思·埃奇沃斯（Kenneth Edgeworth）。$^{[12]}$ 对该结构最为直接的预测来自天文学家胡利奥·安赫尔·费尔南德斯（Julio Ángel Fernández），他在 1980 年发表论文，提出在海王星之外存在一条彗星带$^{[13][14]}$，并认为它可能是短周期彗星的来源。$^{[15][16]}$

1992 年，小行星 15760 Albion 被发现，这是继 1930 年发现冥王星、1978 年发现卡戎之后确认的首个柯伊伯带天体（KBO）。$^{[17]}$ 自其发现以来，已知的柯伊伯带天体数量已增加到数千个，并且人们认为直径超过 100 km（62 英里）的 KBO 可能多达 100000 个以上。$^{[18]}$ 起初，柯伊伯带被认为是周期彗星（轨道周期小于 200 年）的主要储存库。然而，自 20 世纪 90 年代中期以来的研究表明，柯伊伯带在动力学上是稳定的，而彗星真正的起源地是散射盘——这是一个由海王星在约 45 亿年前向外迁移所形成的动力学活跃区域。$^{[19]}$ 散射盘天体（例如阋神星 Eris）具有极高的轨道偏心率，其轨道最远可延伸至距离太阳约 100 AU 的位置。$^{[a]}$

柯伊伯带不同于假想中的奥尔特云，后者被认为距离太阳远得多（约为其一千倍），且整体上近似球形。柯伊伯带内的天体，与散射盘成员以及任何潜在的希尔斯云或奥尔特云天体，统称为海王星外天体（TNO）。$^{[22]}$ 冥王星是目前已知柯伊伯带中最大、质量最大的成员，同时也是已知 TNO 中体积最大、质量第二大的天体，仅次于散射盘中的阋神星。$^{[a]}$ 冥王星最初被视为一颗行星，但由于其属于柯伊伯带成员，这一事实促使其在 2006 年被重新归类为矮行星。冥王星在成分上与许多其他柯伊伯带天体相似，其轨道周期也具有一类 KBO 的典型特征，这类天体被称为“冥族天体”（plutinos），它们与海王星处于相同的 2:3 轨道共振状态。

柯伊伯带与海王星有时也被视为太阳系范围的标志之一，其他替代界限还包括日球层顶（heliopause），以及太阳引力影响与其他恒星引力影响相当的距离（估计约在 50000–125000 AU 之间）。$^{[23]}$
\subsection{历史}
\begin{figure}[ht]
\centering
\includegraphics[width=6cm]{./figures/c3c318a9ff09ac52.png}
\caption{冥王星与卡戎} \label{fig_Kuiper_1}
\end{figure}
1930 年冥王星被发现之后，许多人开始推测它可能并非孤立存在。如今被称为柯伊伯带的这一区域，在随后数十年间以多种不同形式被提出为假想结构。直到 1992 年，人们才首次获得其存在的直接证据。由于在此之前关于柯伊伯带性质的设想在数量和形式上都十分多样，至今仍难以确定究竟应将最早提出这一概念的功劳归于谁。$^{[24]:106}$
\subsubsection{假说}
\begin{figure}[ht]
\centering
\includegraphics[width=6cm]{./figures/acd1d13b7d423c69.png}
\caption{以其命名柯伊伯带的天文学家杰拉德·柯伊伯} \label{fig_Kuiper_2}
\end{figure}